\section*{الفصل \ref{ch:g8:reproductive-systems} --- الجهازان التناسليان}

\begin{solution}{exo:g8:reproductive-systems:1}
\omterm{def:g8:reproductive-systems:male}{الخصيتان}، في كيس الصفن: إنتاج مستمر للنطف منذ \omterm{def:g8:puberty:puberty}{البلوغ}.
القنوات والغدد: التخزين والنقل والسوائل التي تصنع المني.
\omterm{def:g8:reproductive-systems:male}{القضيب}: مخرج المني --- ومخرج \omterm{def:g7:removing-wastes:kidneys}{البول} أيضًا، في مهمته
المنفصلة التي تحرسها \omterm{def:g7:heart-and-circulation:rooms}{الصمامات}.
\end{solution}

\begin{solution}{exo:g8:reproductive-systems:2}
\omterm{def:g8:reproductive-systems:female}{المبيضان}: مخزون \omterm{def:g4:animal-reproduction:sexual}{البويضات} منذ \omterm{def:g2:born-grow-die:cycle}{الولادة}، وإطلاق واحدة في كل
دورة منذ \omterm{def:g8:puberty:puberty}{البلوغ}. \omterm{def:g8:reproductive-systems:female}{قناتا البيض}: التقاط \omterm{def:g4:animal-reproduction:sexual}{البويضة} المنطلقة ---
وهما موضع \omterm{def:g5:flower-to-fruit:steps}{الإخصاب} المعتاد. \omterm{prop:g5:human-reproduction:plan}{الرحم}: المسكن العضلي المتجدد
البطانة. المهبل: استقبال المني، وقناة \omterm{def:g2:born-grow-die:cycle}{الولادة}.
\end{solution}

\begin{solution}{exo:g8:reproductive-systems:3}
\omterm{def:g8:reproductive-systems:male}{النطفة}: أصغر \omterm{def:g5:first-look-at-cells:cell}{خلايا} الجسم --- رأس محشو \omterm{def:g6:cell-unit-of-life:parts}{بالنواة} وذيل
للسباحة، زاده خفيف، ويصنع بلا انقطاع بمئات الملايين.
\omterm{def:g4:animal-reproduction:sexual}{البويضة}: أكبر \omterm{def:g5:first-look-at-cells:cell}{خلايا} الجسم، مستديرة ساكنة، غنية بالزاد،
وتطلق واحدة في الشهر.
\end{solution}

\begin{solution}{exo:g8:reproductive-systems:4}
يطلق أحد \omterm{def:g8:reproductive-systems:female}{المبيضين} \omterm{def:g4:animal-reproduction:sexual}{بويضة} واحدة، تلتقطها \omterm{def:g8:reproductive-systems:female}{قناة البيض} ---
نحو اليوم 14. وتبقى قابلة \omterm{def:g4:animal-reproduction:sexual}{للإخصاب} نحو يوم واحد.
\end{solution}

\begin{solution}{exo:g8:reproductive-systems:5}
تفكك بطانة \omterm{prop:g5:human-reproduction:plan}{الرحم} المهيأة وتساقطها، وخروجها في أيام
قليلة. ووصوله يقول إن تلك الدورة انتهت من غير \omterm{def:g4:animal-reproduction:sexual}{إخصاب}.
\end{solution}

\begin{solution}{exo:g8:reproductive-systems:6}
\omterm{def:g5:flower-to-fruit:steps}{الإخصاب}. تستقر \omterm{prop:g8:sexual-reproduction:core}{اللاقحة} في البطانة المهيأة، فتتوقف
الدورة، ويبدأ \omterm{def:g5:human-reproduction:pregnancy}{الحمل}.
\end{solution}

\begin{solution}{exo:g8:reproductive-systems:7}
الخطوة 1، \omterm{def:g8:sexual-reproduction:gametes}{الأمشاج}: \omterm{def:g8:reproductive-systems:male}{الخصيتان} \omterm{def:g8:reproductive-systems:female}{والمبيضان}. الخطوة 2،
\omterm{def:g5:flower-to-fruit:steps}{الإخصاب}: في \omterm{def:g8:reproductive-systems:female}{قناة البيض}، تبلغها \omterm{def:g8:reproductive-systems:male}{النطفة} سباحةً من المهبل
عبر \omterm{prop:g5:human-reproduction:plan}{الرحم}. الخطوة 3، نمو \omterm{prop:g8:sexual-reproduction:core}{اللاقحة} المنقسمة: في البطانة
المهيأة داخل \omterm{prop:g5:human-reproduction:plan}{الرحم}.
\end{solution}

\begin{solution}{exo:g8:reproductive-systems:8}
يجري الجانب الذكري على منطق الكثرة والرخص: ملايين سابحة
بلا انقطاع، كل واحدة يمكن الاستغناء عنها. ويجري الجانب
الأنثوي على منطق القلة والتزويد: \omterm{def:g4:animal-reproduction:sexual}{بويضة} واحدة مزوّدة في
الشهر، وخلفها استثمار \omterm{prop:g5:human-reproduction:plan}{الرحم} كله في المسكن. إنه خط
الاستراتيجيتين، مكتوبًا في \omterm{def:g5:first-look-at-cells:cell}{الخلايا}.
\end{solution}

\begin{solution}{exo:g8:reproductive-systems:9}
الآلة تضبط نفسها في سنواتها الأولى --- فسلسلة \omterm{def:g8:puberty:hormones}{الهرمونات}
تستقر على إيقاعها --- فعدم الانتظام متوقع؛ وطول الدورة
يختلف فعلًا من شخص إلى آخر ومن شهر إلى آخر. وكلاهما
الجهاز يعمل عملًا عاديًا.
\end{solution}

\begin{solution}{exo:g8:reproductive-systems:10}
\omterm{prop:g5:human-reproduction:plan}{الرحم} \omterm{def:g3:bones-and-muscles:muscle}{عضلة}، والتساقط عمل عضلي --- فتحس انقباضاته
تقلصات. والدفء والحركة يساعدان (ولدى ممرضة المدرسة
نصيحة حين لا يكفيان).
\end{solution}

\begin{solution}{exo:g8:reproductive-systems:11}
مهمة \omterm{def:g8:reproductive-systems:male}{النطفة} هي الرحلة: رخيصة، سابحة، ومطلوبة بالملايين
لأن أحدًا منها لا يكاد يصل --- فيجري مصنعها بلا انقطاع.
ومهمة \omterm{def:g4:animal-reproduction:sexual}{البويضة} أن تصير زاد \omterm{prop:g8:sexual-reproduction:core}{اللاقحة} الأول: مكلفة، مخزونة
منذ \omterm{def:g2:born-grow-die:cycle}{الولادة}، وتطلق وخلفها تهيئة شهر كامل --- فمصنعها
خزانة تنفق واحدة في كل مرة.
\end{solution}

\begin{solution}{exo:g8:reproductive-systems:12}
الحدث 1 هو تهيئة السؤال: البطانة تبنى، والمسكن يهيأ.
والحدث 2 يطرحه: \omterm{def:g4:animal-reproduction:sexual}{بويضة} واحدة تنطلق --- أيأتي \omterm{def:g5:flower-to-fruit:steps}{الإخصاب}؟
والحدث 3 هو الجواب ``لا'': تتساقط البطانة التي لم تعد
لازمة، وتسأل العجلة من جديد في الشهر المقبل. والحدث 4 هو
الجواب ``نعم'': تستقر \omterm{prop:g8:sexual-reproduction:core}{اللاقحة}، ويتوقف السؤال، ويؤدي
المسكن غرضه تسعة أشهر.
\end{solution}

\begin{solution}{pb:g8:reproductive-systems:1}
\textbf{1.} من \omterm{def:g8:reproductive-systems:male}{الخصيتين} --- المصنعين المحفوظين في كيس
الصفن --- تمر \omterm{def:g8:reproductive-systems:male}{النطف} إلى القنوات للتخزين والنقل، وتنضم
إليها سوائل الغدد فيتكوّن المني، ثم تخرج عبر \omterm{def:g8:reproductive-systems:male}{القضيب}؛
والعضو نفسه يحمل \omterm{def:g7:removing-wastes:kidneys}{البول} في أوقات أخرى، وترتيب من الصمامات
يفصل بين المهمتين فصلًا تامًا.

\textbf{2.} ينتظر المخزون في \omterm{def:g8:reproductive-systems:female}{المبيضين}، منذ \omterm{def:g2:born-grow-die:cycle}{الولادة}؛ ويحدث \omterm{def:g5:flower-to-fruit:steps}{الإخصاب}
عادة في \omterm{def:g8:reproductive-systems:female}{قناتي البيض}؛ ويستقر \omterm{def:g5:human-reproduction:pregnancy}{الحمل} في بطانة \omterm{prop:g5:human-reproduction:plan}{الرحم}
المهيأة.

\textbf{3.} المشيجان مبنيان على استراتيجيتين متقابلتين: ملايين من
السابحات الرخيصة في مقابل \omterm{def:g4:animal-reproduction:sexual}{بويضة} واحدة مزوّدة في الشهر ---
موازنة العدد والرعاية، مكتوبة في \omterm{def:g5:first-look-at-cells:cell}{الخلايا} نفسها.

\textbf{4.} كانا ينتظران أوامرهما: قام الجهازان مبنيين لكن ساكنين
حتى شغّلتهما هرمونات مركز التحكم في قاعدة \omterm{def:g5:brain-in-command:system}{الدماغ} عند
\omterm{def:g8:puberty:puberty}{البلوغ}.

\textbf{5.} اليوم 1 إلى 5 تقريبًا: \omterm{prop:g8:reproductive-systems:cycle}{الحيض}. وحتى منتصف الدورة: تبنى
البطانة. ونحو اليوم 14: \omterm{prop:g8:reproductive-systems:cycle}{الإباضة}، ومعها اليوم الخصب.
ثم الانتظار المهيأ --- إلى اليوم 28 ودورة العجلة
التالية.

\textbf{6.} يختلف طول الدورة من شخص إلى آخر، ومن شهر إلى آخر عند
الشخص الواحد --- والسنوات الأولى تجري غير منتظمة بينما
تضبط الآلة نفسها. فالتقويم يعلّم النمط، لا تواريخ أحد
بعينه.

\textbf{7.} الملغى: تفكك البطانة \omterm{prop:g8:reproductive-systems:cycle}{والحيض} التالي --- فالعجلة تقف.
والبادئ: استقرار \omterm{prop:g8:sexual-reproduction:core}{اللاقحة} في البطانة، وأشهر \omterm{def:g5:human-reproduction:pregnancy}{الحمل}
التسعة.

\textbf{8.} التقلصات هي \omterm{prop:g5:human-reproduction:plan}{الرحم} العضلي يعمل وهو يتساقط --- أمر عادي،
لا ضرر؛ والدفء والحركة يساعدان، ولما هو أشد من ذلك
عنوانه عند باب الممرضة أو الطبيب.

\textbf{9.} مثلًا: ``هذه أعضاء \omterm{def:g4:heart-and-blood:heart}{كالقلب} \omterm{def:g7:how-animals-breathe:four}{والرئتين} --- وكل من في هذه
القاعة يملك إحدى المجموعتين، وهذا ببساطة فصلها.''

\textbf{10.} المصانع: \omterm{def:g8:reproductive-systems:male}{الخصيتان} --- \omterm{def:g8:reproductive-systems:female}{المبيضان}. \omterm{def:g8:sexual-reproduction:gametes}{الأمشاج}: ملايين سابحة
بلا انقطاع --- \omterm{def:g4:animal-reproduction:sexual}{بويضة} واحدة مزوّدة في الشهر. اللقاء:
\omterm{def:g8:reproductive-systems:female}{قناة البيض}. المسكن: لا مسكن --- \omterm{prop:g5:human-reproduction:plan}{والرحم} المهيأ شهريًا.
الجدول الزمني: مستمر منذ \omterm{def:g8:puberty:puberty}{البلوغ} --- ودوري، من \omterm{def:g8:puberty:puberty}{البلوغ} إلى
نحو الخمسين.

\textbf{11.} الخطوة 1 على النموذجين: المصنعان يصنعان
\omterm{def:g8:sexual-reproduction:gametes}{الأمشاج}. والخطوة 2 على النموذج الأنثوي: التقاؤهما في
\omterm{def:g8:reproductive-systems:female}{قناة البيض}. والخطوة 3 عبر النموذج والتقويم: استقرار
\omterm{prop:g8:sexual-reproduction:core}{اللاقحة} في بطانة التقويم المهيأة، والعجلة واقفة، والنمو
قد بدأ.

\textbf{12.} ``لأن الأجساد تسير على المعرفة خيرًا مما
تسير على الأسطورة --- ولكل ما عُلّم اليوم عنوان مدرَّب
سري تُتابع عنده أسئلته.''
\end{solution}
